\documentclass[a4paper,12pt,fleqn]{article}
\usepackage[ngerman]{babel}
\usepackage[utf8]{inputenc}
\usepackage[T1]{fontenc}
\usepackage{amsmath,amssymb}
\usepackage{tikz}
\usetikzlibrary{automata,positioning,arrows.meta}

\tikzset{
    ->,>=Stealth,
    node distance=2.5cm,
    every state/.style={thick,minimum size=1cm},
    initial text={},
}

\newcommand{\methodbox}[1]{%
\medskip
\noindent\fbox{%
    \begin{minipage}{0.94\linewidth}
    \textbf{Methodik / Definition.}\par\smallskip
    \small
    #1
    \end{minipage}%
}
\medskip\par
}

\newcommand{\examplebox}[1]{%
\medskip
\noindent\fbox{%
    \begin{minipage}{0.94\linewidth}
    \textbf{Beispiele.}\par\smallskip
    \small
    #1
    \end{minipage}%
}
\medskip\par
}

\newcommand{\taskbox}[1]{%
\medskip
\noindent\fbox{%
    \begin{minipage}{0.94\linewidth}
    \textbf{Aufgabenstellung.}\par\smallskip
    \small
    #1
    \end{minipage}%
}
\medskip\par
}

\begin{document}
\raggedright
\setlength{\parindent}{0pt}
\setlength{\mathindent}{0pt}

{\LARGE\bfseries ETI -- Aufgabenblatt 6 \\ Praesenzaufgaben\par}
\vspace{2em}

\section*{Aufgabe 6.1}

\taskbox{
Entwerfen Sie einen PDA mit hoechstens $4$ Zustaenden und hoechstens $2$ Zeichen im
Kelleralphabet, der die Sprache
\[
    Q=\{w\in\{0,1\}^*\mid |w|_0\ne |w|_1\}
\]
erkennt. Begruenden Sie kurz die Korrektheit. Ist der Automat deterministisch?
}

\methodbox{
Der Zustand speichert das Vorzeichen der Differenz: In $q_0$ wurden bisher mehr $0$en
gelesen, in $q_1$ bisher mehr $1$en, und in $q_=$ ist die Differenz $0$.
Damit nur zwei Kellerzeichen benoetigt werden, speichern wir im Keller nicht die ganze
Differenz, sondern $|\ |w|_0-|w|_1\ |-1$ viele $A$-Symbole ueber dem Bodensymbol $\$$.
Eine Differenz vom Betrag $1$ wird also nur durch den Zustand und das Bodensymbol dargestellt.
}

Wir verwenden
\[
    \Gamma=\{A,\$\}
\]
und akzeptieren per Endzustand. Der PDA ist
\[
P=(\{q_s,q_=,q_0,q_1\},\{0,1\},\{A,\$\},\delta,q_s,\{q_0,q_1\}).
\]
Dabei ist $q_s$ nur der Startzustand, der das Bodensymbol in den Keller legt.

\begin{center}
\begin{tikzpicture}[node distance=3.1cm and 3.7cm]
    \node[state, initial] (qs) {$q_s$};
    \node[state, right=of qs] (qe) {$q_=$};
    \node[state, accepting, above right=of qe] (q0) {$q_0$};
    \node[state, accepting, below right=of qe] (q1) {$q_1$};

    \path (qs) edge node[above] {$\varepsilon,\varepsilon\to\$$} (qe)
          (qe) edge[bend left=16] node[above,sloped] {$0,\$\to\$$} (q0)
          (q0) edge[bend left=16] node[below,sloped] {$1,\$\to\$$} (qe)
          (qe) edge[bend right=16] node[below,sloped] {$1,\$\to\$$} (q1)
          (q1) edge[bend right=16] node[above,sloped] {$0,\$\to\$$} (qe)
          (q0) edge[loop above] node {$\begin{array}{c}
              0,\$\to A\$\\
              0,A\to AA\\
              1,A\to\varepsilon
          \end{array}$} ()
          (q1) edge[loop below] node {$\begin{array}{c}
              1,\$\to A\$\\
              1,A\to AA\\
              0,A\to\varepsilon
          \end{array}$} ();
\end{tikzpicture}
\end{center}

\examplebox{
Fuer $w=001$ gilt:
\[
    \$ \xrightarrow{0} \$ \xrightarrow{0} A\$ \xrightarrow{1} \$.
\]
Der Automat endet in $q_0$, also wird akzeptiert, denn es gibt eine $0$ mehr als $1$en.

Fuer $w=0101$ wechselt der Automat nach jedem Ausgleich zurueck nach $q_=$:
\[
    q_= \xrightarrow{0} q_0 \xrightarrow{1} q_=
    \xrightarrow{0} q_0 \xrightarrow{1} q_=.
\]
Das Wort wird nicht akzeptiert, weil am Ende gleich viele $0$en und $1$en gelesen wurden.
}

Zur Korrektheit betrachten wir nach jedem gelesenen Praefix $u$ die Differenz
\[
    d(u)=|u|_0-|u|_1.
\]
Ist $d(u)=0$, befindet sich der Automat in $q_=$ und im Keller liegt nur $\$$. Ist
$d(u)>0$, befindet er sich in $q_0$; liegen $k$ viele $A$-Symbole auf dem Keller, dann gilt
$d(u)=k+1$. Ist $d(u)<0$, befindet er sich entsprechend in $q_1$ und bei $k$ vielen
$A$-Symbolen gilt $-d(u)=k+1$.

\medskip
Liest der Automat ein Zeichen, das die aktuelle Differenz vergroessert, bleibt er im
gleichen Vorzeichenzustand und erhoeht den Kellerinhalt passend. Liest er ein Zeichen, das
die Differenz verkleinert, wird ein $A$ geloescht; beim Uebergang von Betrag $1$ zu Betrag
$0$ steht oben bereits $\$$, und der Automat wechselt nach $q_=$. Damit gilt nach dem ganzen
Wort:
\[
    w\in Q
    \quad\Longleftrightarrow\quad
    d(w)\ne 0
    \quad\Longleftrightarrow\quad
    \text{der Automat endet in } q_0 \text{ oder } q_1.
\]
Also erkennt der PDA genau $Q$.

\medskip
Der Automat ist deterministisch. Fuer jede Kombination aus Zustand, naechstem Eingabezeichen
und Kelleroberseite gibt es hoechstens einen Uebergang. Ausser dem einmaligen Startuebergang
aus $q_s$ gibt es keine $\varepsilon$-Uebergaenge, die mit Eingabeuebergaengen konkurrieren.

\end{document}
